\section{Related work}
\label{sec:related}

\paragraph{Prophet and Stan.}
\citet{taylor2018forecasting} introduce the changepoint prior
$\delta_j \sim \mathrm{Laplace}(0, \tau)$ as a sparse prior for selecting
changepoints automatically, and fit the model with ``Stan's L-BFGS to find a
maximum a posteriori estimate''. \citet{carpenter2017stan} describe Stan's
optimizers, L-BFGS by default, with reference to \citet{nocedal2006numerical}, and a
run that ends as soon as any of its tolerances on the log density, the gradient or
the parameters is met. Both descriptions are accurate. This paper is about what
happens where the two meet.

\paragraph{Quasi-Newton methods on nonsmooth objectives.}
The textbook treatment of L-BFGS and Newton's method assumes smoothness
\citep{nocedal2006numerical}. \citet{lewis2013nonsmooth} study BFGS on nonsmooth
functions and find it more capable than that assumption suggests, given a line
search on the weak Wolfe condition; their suggestion that L-BFGS codes offer such a
line search is the nearest thing in that literature to the change
\cref{sec:smooth} points at. \citet{andrew2007scalable} start from the observation
that L-BFGS cannot be used as it stands for an $\ell_1$-regularised loss and modify
it, as OWL-QN, by keeping each step within an orthant. An earlier version of this
implementation used OWL-QN before the reformulation of \cref{sec:reformulation}
replaced it.

\paragraph{Methods built for the $\ell_1$ term.}
Besides orthant-wise quasi-Newton methods, cyclical coordinate descent along a
regularisation path \citep{friedman2010regularization} and iterative
shrinkage-thresholding \citep{beck2009fast} handle an $\ell_1$ penalty directly.
Each of them changes the method. The reformulation of \cref{sec:reformulation}
changes the problem instead: written as a difference of non-negative parts, the
penalty becomes linear and the bounds are left to a bound-constrained quasi-Newton
method \citep[L-BFGS-B;][]{byrd1995limited}. The device is old.
\citet[Section~6]{tibshirani1996regression} gives it, attributing it to David Gay,
as one of two algorithms for the lasso; \citet{chen2001atomic} use it to state basis
pursuit as a linear program, noting that the equivalence of minimum-$\ell_1$
problems with linear programming had been known since the 1950s; and
\citet{figueiredo2007gradient} solve the resulting bound-constrained quadratic
program by gradient projection, which is the closest in spirit to what is done
here. What \cref{sec:reformulation} adds is the statement of exactness for
Prophet's full posterior, including its non-convex configurations, and for the
local minimisers and first-order points that a solver can actually deliver.

\paragraph{The lasso and its solution path.}
With $\sigma$ fixed, linear growth and additive seasonality, Prophet's MAP problem
is a least-squares problem with Gaussian penalties on $k$, $m$ and $\beta$ and an
$\ell_1$ penalty on $\delta$: a lasso-type problem \citep{tibshirani1996regression}.
For the lasso itself the solution path is piecewise linear in the penalty and can
be computed exactly \citep{efron2004least}; the lasso literature therefore offers a
standard against which a sequence of Prophet fits across $\tau$ can be read
(\cref{sec:consequences}).

\paragraph{What is not new.}
None of the methods above is ours, and neither is the observation that L-BFGS is a
poor fit for an $\ell_1$ term. The contribution is the four items listed in
\cref{sec:introduction}: a diagnosis in a specific, widely deployed system, the
controls that isolate it, an exact fix stated for that system's full model, and the
attribution of its user reports.
